\documentclass[]{msu-thesis}

\usepackage[T1]{fontenc}
\usepackage{natbib}

\title{On the Configurability of Long Documents}
\author{Noor Fatima}
\fieldofstudy{Computer Science}
\date{2026}

\begin{document}

\frontmatter
\maketitlepage

\begin{abstract}
This thesis argues that a single, well-structured class can carry the full
weight of graduate thesis production, from title page through bibliography.
We demonstrate the approach with a worked chapter, a cited source, and a
compiled bibliography.
\end{abstract}

\clearpage
\makecopyrightpage

\clearpage
\SingleSpacing
\tableofcontents*

\mainmatter

\chapter{Introduction}
Thesis templates tend to specialise per institution. \citet{knuth1984texbook}
shows that the underlying typesetting engine does not require this — the
constraints come from institutional convention, not from \TeX{} itself.

\section{Structure of this Document}
The remaining chapter extends this argument with a worked example.

\chapter{Related Work}
Prior systems addressing long-document typesetting are surveyed in
\citet{knuth1984texbook}.

\bibliographystyle{plainnat}
\bibliography{references}

\begin{appendix}
\chapter{Supplementary Material}
Additional material supporting the chapters above.
\end{appendix}

\end{document}
